\scanpage{149}\mspage{4}%
% First occurrence of strictes: one wavy underline; second occurrence: one straight underline.
pé at length and in detail in Madame Sinh's book, I think, under the chapter title of ``\uline{strict} Gr-categories'' (i.e. \inserted{\unsure{where}} the associativity, unit, and inverse isomorphisms $X\mathbin{.}X^{-1}\unsure{\simeq|=}1$ are \uline{identities}). One can make explicit for example how (via the fact that every Gr-category is Gr-equivalent to a strict Gr-category) the calculus with Gr-categories \unsure{reduces} to a very algebraic calculus with \uline{strict} Gr-categories, by a sort of ``calculus of fractions'' (on the right or on the left, as one chooses) of the type of the one that leads to the construction of derived categories. Be that as it may, here is the explicit description of this structure (group in $(1\text{-ct.})$) in terms of group theory ($1$-category in $(\mathrm{Gr})$) :

The structure is described by a \unsure{.} quadruple $(L_1,L_0,d,\unsure{\Theta|\theta})$, with
\diagram{\begin{tikzcd}[ampersand replacement=\&]
L_1 \arrow[r,"d"] \& L_0\,\raisebox{1.1ex}{$\scriptstyle\unsure{:}$}
\end{tikzcd}}
a homomorphism of ordinary groups, and
\[
\Theta:L_0\longrightarrow\underline{\mathrm{Aut}}_{\mathrm{gr}}(L_1)
\]
an action of $L_0\deleted{\unsure{\Theta|\theta|0}}$ on $L_1$, with the
% archive sheet number: 149
